% ----------------------
\subsection{Section 53 (8 June 1831)}
% ----------------------
	\label{DC53}
	
	% -----
	\subsubsection{Historical Background}
	% -----
		\label{DC53-HB}
		
		% --
		\paragraph{JSP}
		% --
		
			``Sidney Gilbert, the recipient of this revelation, was Newel K. Whitney's business partner and an early convert to the church. He may have been present at a 5 June 1831 meeting at which JS taught that `the Elders would have large congregations to speak to and they must soon take there departure into the Reagions west.'1 The next evening, JS dictated a revelation that commissioned a number of elders to accompany him to Missouri, which the text called `the land which I will consecrate unto my People which are a remnant of Jacob \& those who are heirs according to the covenant.' The 6 June revelation did not name Gilbert among those who should travel to Missouri. JS's history records that shortly after the reception of that revelation, Gilbert approached JS and requested direction for himself. In response, JS dictated this revelation that gave Gilbert assignments in the church, including the instruction to travel to Missouri with JS and Sidney Rigdon. Gilbert left for Missouri on 18 June with JS and Rigdon.''
		
	% -----
	\subsubsection{Versions}
	% -----
		\label{DC53-V}
		
		See the \href{https://drive.google.com/file/d/1goAvNz8jS4R8slYfMG_db55Ty2z_vmgK/view?usp=sharing}{sections comparison} document.
	
		\begin{compactdesc}
			\item[JSP] \hyperref[EMS-1831-JS-8-JUN-1831]{8 June 1831}
			\item[BC] Chapter 55, 127--128
			\item[DC 1835] Section 66,\footnotemark 195
			\item[TS] 5:4 (15 February 1844), 432
			\item[DC 1844] Section 67, 296--297
			\item[MS] 5:8 (January 1845), 113
		\end{compactdesc}
			\footnotetext{%
				% \label{}%
				As HDDC 695 notes, the number of section 66 in the 1835 DC is a mistake, since it should be 67.
				}

	% -----
	\subsubsection{Section 53:1} % *DC53-1 
	% -----
		\label{DC53-1}

		BEHOLD, I say unto you, my servant Sidney Gilbert, that I have heard your prayers; and you have called upon me that it should be made known unto you, of the Lord your God, concerning your calling and election in the church, which I, the Lord, have raised up in these last days.

		\begin{framed}
		\scriptsize
			\begin{compactdesc}
				\item[JSP] Behold I say unto you my servent Sidney that I have heard your prayers \& ye have called upon me that it should be made known unto you of the lord your God concerning your calling \& election in this Church which I the Lord have raised up in these last days
				% \item[BC] 
				% \item[]
			\end{compactdesc}
		\end{framed}

	% -----
	\subsubsection{Section 53:2} % *DC53-2 
	% -----
		\label{DC53-2}

		Behold, I, the Lord, who was crucified for the sins of the world, give unto you a commandment that you shall forsake the world.

		\begin{framed}
		\scriptsize
			\begin{compactdesc}
				\item[JSP] Behold I \sout{my} the Lord who was crusified for the sins of the world giveth unto you a commandment that you shall forsake the world
				% \item[BC] 
				% \item[]
			\end{compactdesc}
		\end{framed}

	% -----
	\subsubsection{Section 53:3} % *DC53-3 
	% -----
		\label{DC53-3}

		Take upon you mine \hyperref[ORDINATION]{ordination},%
			\footnote{%
				% \label{}%
				Notice that all the older versions have ``\hyperref[ORDINANCE]{ordinances}'' here. This is a potentially significant change in thinking about the \hyperref[PRIESTHOOD]{priesthood} and how the general understanding of it has changed. See \hyperref[DC53-6]{verse 6}, which makes more sense in the light of the older versions of this verse. See also OED definition \#I.5 for ``ordinance,'' ``An authoritative instruction as to how to proceed or act; an established set of principles; a system of government; authority; discipline.'' This definition makes sense of ``ordinance'' in both verses here. There is also OED \#I.4.a, ``A practice or usage authoritatively enjoined or prescribed; esp. a religious or ceremonial observance, as the sacraments, etc.''
				}
		even that of an elder, to preach faith and repentance and remission of sins, according to my word, and the reception of the Holy Spirit by the laying on of hands;

		\begin{framed}
		\scriptsize
			\begin{compactdesc}
				\item[JSP] take upon you mine ordinances even that of an Elder to Preach faith \& repentance \& remission of sins according to my word \& the reception of the Holy spirit by the laying on of hands
				\item[BC] 3 Take upon you mine ordinances, even that of an elder, to preach faith and repentance, and remission of sins, according to my word, and the reception of the Holy Spirit by the laying on of hands.
				\item[DC 1835] Take upon you mine ordinances, even that of an elder, to preach faith and repentance, and remission of sins, according to my word, and the reception of the Holy Spirit by the laying on of hands.
				\item[DC 1844] Take upon you mine ordinances, even that of an elder, to preach faith and repentance, and remission of sins, according to my word, and the reception of the Holy Spirit by the laying on of hands.
			\end{compactdesc}
		\end{framed}

	% -----
	\subsubsection{Section 53:4} % *DC53-4 
	% -----
		\label{DC53-4}

		And also to be an agent unto this church in the place which shall be appointed by the bishop, according to commandments which shall be given hereafter.

		% \begin{framed}
		% \scriptsize
			% \begin{compactdesc}
				% \item[JSP] 
				% % \item[BC] 
				% % \item[]
			% \end{compactdesc}
		% \end{framed}

	% -----
	\subsubsection{Section 53:5} % *DC53-5 
	% -----
		\label{DC53-5}

		And again, verily I say unto you, you shall take your journey with my servants Joseph Smith, Jun., and Sidney Rigdon.

		% \begin{framed}
		% \scriptsize
			% \begin{compactdesc}
				% \item[JSP] 
				% % \item[BC] 
				% % \item[]
			% \end{compactdesc}
		% \end{framed}

	% -----
	\subsubsection{Section 53:6} % *DC53-6 
	% -----
		\label{DC53-6}

		Behold, these are the first \hyperref[ORDINANCE]{ordinances}%
			\footnote{%
				% \label{}%
				Note the change from ``ordinances'' to ``ordination'' in different versions of \hyperref[DC53-3]{verse 3}.
				}
		which you shall receive; and the residue shall be made known in a time to come,%
			\footnote{%
				% \label{}%
				Possibly looking forward to temple ordinances.
				}
		according to your labor in my vineyard.

		% \begin{framed}
		% \scriptsize
			% \begin{compactdesc}
				% \item[JSP] 
				% % \item[BC] 
				% % \item[]
			% \end{compactdesc}
		% \end{framed}

	% -----
	\subsubsection{Section 53:7} % *DC53-7 
	% -----
		\label{DC53-7}

		And again, I would that ye should learn that he only is saved who endureth unto the end.  Even so.  Amen.

		% \begin{framed}
		% \scriptsize
			% \begin{compactdesc}
				% \item[JSP] 
				% % \item[BC] 
				% % \item[]
			% \end{compactdesc}
		% \end{framed}